% Content for the Connecting Technology to Emotions brochure.
\colorlet{Accent}{Violet}
\WorkshopHeader
  {03}
  {CONNECTING TECHNOLOGY\\TO EMOTIONS}
  {Build campus capability in human-centred AI, responsible technology and experience innovation}
  {Human insight \textbullet\ Responsible AI \textbullet\ Interdisciplinary innovation}
  {\faHeart}
  {assets/bg-technology-emotions.png}

\MetaStrip
  {Students, teachers \& designers}
  {1 day \textbar\ 9:30--17:00}
  {Design \& prototype lab}
  {Discuss with us}

\begin{minipage}[t]{0.615\textwidth}
  \vspace{0pt}
  \begin{brochurecard}[colback=Accent!3!white]
    \SectionHead{Why this matters now}
    \LeaderHook{Position the institution at the intersection of human-centred AI, responsible technology and interdisciplinary innovation.}
    \WhyItMatters
      {Students learn to design around trust, emotion, inclusion and ethics---capabilities central to AI, products and meaningful technology adoption.}
      {Academicians open interdisciplinary research directions, differentiated IP and funding opportunities in human-centred AI and responsible technology.}
  \end{brochurecard}

  \begin{brochurecard}
    \begin{minipage}[t]{0.485\linewidth}
      \SectionHead{What we will explore}
      \begin{itemize}
        \item Recognise emotional tensions as sources of novel technology problems.
        \item Use stories and empathy evidence to guide responsible solution design.
        \item Balance usefulness, privacy, trust and inclusion through inventive frameworks.
      \end{itemize}
    \end{minipage}\hfill
    \begin{minipage}[t]{0.485\linewidth}
      \SectionHead{What participants gain}
      \begin{itemize}
        \item \textbf{Find unmet needs:} reveal novel human--technology problems worth researching.
        \item \textbf{Deepen contribution:} move beyond feature iterations to meaningful experience innovation.
        \item \textbf{Create defensible ideas:} shape non-obvious concepts with patent potential.
        \item \textbf{Build support:} present an ethical, partner-ready case for collaboration and funding.
      \end{itemize}
    \end{minipage}
  \end{brochurecard}

  \begin{brochurecard}
    \SectionHead{One-day learning journey}
    \ScheduleRow{09:30--10:15}{\textbf{Technology is an emotional experience:} trust, anxiety, agency and belonging as innovation signals.}
    \ScheduleRow{10:15--11:15}{\textbf{Find unmet needs:} empathy interviews, emotional vocabulary and bias-aware observation.}
    \ScheduleRow{11:30--12:45}{\textbf{Map the experience:} personas, emotion journeys, moments that matter and design principles.}
    \ScheduleRow{13:45--15:00}{\textbf{Invent around tensions:} use TRIZ for privacy--personalisation, ease--control and speed--care.}
    \ScheduleRow{15:15--16:15}{\textbf{Reverse brainstorm and IP lens:} transform failure modes into differentiated solution concepts.}
    \ScheduleRow{16:15--17:00}{\textbf{Prototype with care:} pitch value, patent potential, emotional signals and ethical safeguards.}
  \end{brochurecard}
\end{minipage}\hfill
\begin{minipage}[t]{0.36\textwidth}
  \vspace{0pt}
  \begin{brochurecard}[colback=Accent!5!white]
    \SectionHead{Who should attend}
    \begin{itemize}
      \item Engineering, design, management and psychology students
      \item Teachers and researchers in technology-enabled learning, HCI and AI
      \item Product, UX and innovation teams in academic incubators
    \end{itemize}
  \end{brochurecard}

  \begin{brochurecard}
    \SectionHead{Methods \& tools}
    \begin{itemize}
      \item Emotion-led problem discovery
      \item Empathy interviews and journey maps
      \item Human-centred and inclusive design
      \item Foundations of affective computing
      \item TRIZ, reverse brainstorming and IP lens
      \item Partner pitch, trust and responsible testing
    \end{itemize}
  \end{brochurecard}

  \begin{brochurecard}[colback=Accent!3!white]
    \SectionHead{Takeaway toolkit}
    \begin{itemize}
      \item An emotion journey map grounded in empathy evidence
      \item Human--technology tensions translated into design principles
      \item Differentiated concepts generated through TRIZ and reverse brainstorming
      \item A low-fidelity prototype with trust, inclusion and ethics checks
    \end{itemize}
    {\color{Accent}\bfseries\fontsize{6.6}{7.2}\selectfont INSTITUTIONAL RETURN}\par
    New interdisciplinary projects, responsible-AI research directions and partner-ready technology concepts.
  \end{brochurecard}

  \begin{customprogramcard}
    \SectionHead{Custom program}
    {\bfseries Custom programs can be discussed.}\par
    The scope can be shaped around the cohort, institutional priorities and desired outcomes.
    Workbooks, empathy/IP/ethics canvases, prototype materials and participation certificates can be included.
  \end{customprogramcard}
\end{minipage}

\Footer
